%% THERMODYNAMIQUE %%
%
\section{Humidit� relative � partir de la temp�rature du point
de ros�e, seuill�e � 101\%}

\itemstyle{ Op�rateur~: } \operfmt{HumW\_T\_Td}
\index{HumW\_T\_Td}

%%%%%OBJET%%%%%
\itemstyle{ Objet~: }
Cet op�rateur calcule l'humidit� relative � l'eau liquide � partir
des temp�ratures de l'air et du point de ros�e.

%%%%%ARGUMENTS%%%%%
\bigskip
%%%%%ENTREES%%%%
\begin{tabarg}
\multicolumn{3}{l}{\titretab{ Entr\'ees}}\\
\hline
$T_s$ & Vecteur, \texttt{float} & Temp\'erature statique de l'air ($K$ ou $^\circ C$) \\
$T_d$ & Vecteur, \texttt{float} & Temp\'erature du point de ros\'ee ($K$ ou $^\circ C$) 
\end{tabarg}

\medskip
%%%%%SORTIES%%%%
\begin{tabarg}
\multicolumn{3}{l}{\titretab{Sorties}}\\
\hline
$H_u$ & Vecteur, \texttt{float} & Humidit\'e relative \`a l'eau liquide ($\%$)\\
\end{tabarg}


%%%%%FORMULE%%%%
\itemstyle{ Formule }

%%%% THERMODYNAMIQUE
\[
H_u = 100 \times \frac{ e_w(T_d)}{e_w(T_s)} 
\]
Si $T_d > T_s$ alors $H_u$ est forc\'e \`a 101\%.



%%%%COMMENTAIRE%%%%
\itemstyle{Commentaire~:}

La fonction $ew()$ est d�crite en annexe
\ref{annexe-thermo}.

Les formules employ�es correspondent � l'hypoth�se que 
l'argument d'entr�e $T_d$ est toujours un point de ros�e (et non
un point de givre si $T_d < 0$�C); 
si $T_d$ provient d'un calcul, cette hypoth�se peut �tre v�rifi�e,
mais si il provient d'un hygrom�tre � miroir
refroidi, il est probable que les valeurs en-dessous de 0�C soient
une mesure de point de givre, et il vaut mieux utiliser l'op�rateur
\hyperref[OP-HumW_T_Tdf]{\operfmt{HumW\_T\_Tdf}}.


%%%%%AUTEUR%%%%%
\itemstyle{ Auteur~:}
CNRM/GMEI/TRAMM

%%%%%REFERENCES%%%%%
\itemstyle{ R�f�rences~:}

\emph{M�t�orologie G�n�rale}, Triplet-Roche, �dition de 1971\cite{Triplet-Roche}, p59.
